\documentclass[11pt]{article}
\usepackage[margin=0.9in]{geometry}
\usepackage{amsmath,amssymb,amsthm,booktabs}
\usepackage[T1]{fontenc}
\usepackage{lmodern}
\usepackage{hyperref}
\hypersetup{colorlinks=true,urlcolor=blue,citecolor=blue,linkcolor=blue,
pdftitle={OEIS A378495, FORMULA (revision 70): refutation and correction},
pdfauthor={Haobo Ma and Wenlin Zhang}}
\newtheorem{theorem}{Theorem}
\newtheorem{lemma}[theorem]{Lemma}
\newtheorem{corollary}[theorem]{Corollary}
\newcommand{\Der}{D}
\title{OEIS A378495, FORMULA (revision \#70):\\refutation and correction}
\author{Haobo Ma\\\small ChronoAI Pte Ltd, Singapore
\and Wenlin Zhang\\\small National University of Singapore\\\small ChronoAI Pte Ltd, Singapore}
\date{November 2, 2026}
\begin{document}
\maketitle
\begin{abstract}
We refute the adjacent-column \emph{Conjecture} and prove the terminal-column
\emph{Conjecture} in the FORMULA section of Peter Kagey's OEIS entry A378495,
revision \#70 of October 17, 2026~\cite{triangle}.
For derangements without a prescribed cycle length, the corrected adjacent
difference is $n!/[2(n-2)]$ for $n\geq5$; at $n=4$ it is $3$.
The literal conjecture fails at $n=5$, where its two sides are $20$ and $6$.
One long-cycle count gives both difference identities and explains the exceptional
boundary. The earliest proof found of these identities, the index correction,
and the companion row-gcd theorem is Su's October 18, 2026 preprint
\cite[Theorem~7.3, Proposition~7.4, Remark~7.5]{su}; Kallat gives a later
independent proof on October 29~\cite[Theorem~4.1]{valuation}.
This note gives a self-contained treatment of the differences and their boundary.
\end{abstract}

\section{The question and the corrected formulas}

Let $\Der_n$ be the number of derangements of $\{1,\ldots,n\}$, with
$\Der_0=1$ and $\Der_1=0$. For $2\leq k\leq n$, let $T(n,k)$ count those
derangements having no cycle of length $k$. We use the displayed triangle's
convention $T(n,1)=0$~\cite{triangle}. This first column is a convention:
literal avoidance of $1$-cycles would instead give $\Der_n$.
Every cycle-avoidance assertion below concerns $k\geq2$.

Write
\begin{equation}\label{eq:b}
 b(m)=(-1)^m\frac{(m+3)!}{2(m+1)}\qquad(m\geq0).
\end{equation}
This is A238474 with its published offset $0$~\cite{bsequence}.
The two FORMULA conjectures in A378495 are
\begin{align}
 T(n,n-1)-T(n,n-2)&=|b(n-4)| &&(n\geq4),\label{eq:original}\\
 T(n,n-2)-T(n,n)&=\frac{(n-3)!(n-4)(n-1)}2 &&(n\geq5).
 \label{eq:terminal}
\end{align}
The first formula needs both an index correction and a separate boundary value.

\begin{theorem}\label{thm:main}
For $n\geq5$,
\begin{align}
 T(n,n-1)-T(n,n-2)
 &=\frac{n!}{2(n-2)}=|b(n-3)|,\label{eq:corrected}\\
 T(n,n-2)-T(n,n)
 &=\frac{(n-3)!(n-4)(n-1)}2.\label{eq:proved}
\end{align}
At $n=4$, both differences on the left are $3$.
Equation~\eqref{eq:original} holds at $n=4$ and fails for every $n\geq5$.
\end{theorem}

Su's Proposition~7.4~\cite{su} proves both identities and records the
$n-3$ correction, the $n=5$ counterexample, and the separate $n=4$ boundary;
Remark~7.5 distinguishes the corrected formula from the printed conjecture.
Kallat's later proof appears in
\cite[Theorem~4.1(c)--(d), Appendix~A.11.3, footnote~3]{valuation},
and his Section~1.5 cites Su's earlier work.
The proof below uses their direct cycle decomposition.

\section{One cycle count and its consequences}

\begin{lemma}\label{lem:long}
If $2\leq k\leq n$ and $2k>n$, then
\begin{equation}\label{eq:long}
 \Der_n-T(n,k)=\binom nk(k-1)!\Der_{n-k}
             =\frac{n!}{k(n-k)!}\Der_{n-k}.
\end{equation}
\end{lemma}
\begin{proof}
A permutation of $n$ points contains at most one $k$-cycle when $2k>n$.
Choose its $k$ labels in $\binom nk$ ways, give them a cyclic order in
$(k-1)!$ ways, and derange the other $n-k$ labels in $\Der_{n-k}$ ways.
The distinguished cycle is unique, so this construction counts every excluded
derangement exactly once. Subtracting the excluded count from $\Der_n$
gives~\eqref{eq:long}.
\end{proof}

\begin{proof}[Proof of Theorem~\ref{thm:main}]
An $n$-cycle has $(n-1)!$ cyclic orders. An $(n-1)$-cycle would leave one
fixed point, so it cannot occur in a derangement. For $n\geq5$, an $(n-2)$-cycle
leaves exactly two labels, which must form a transposition. Thus
\begin{align}
 T(n,n)&=\Der_n-(n-1)!,\label{eq:last}\\
 T(n,n-1)&=\Der_n,\label{eq:next}\\
 T(n,n-2)&=\Der_n-\binom n2(n-3)!.\label{eq:third}
\end{align}
These are also the cases $k=n,n-1,n-2$ of Lemma~\ref{lem:long}, using
$\Der_0=1$, $\Der_1=0$, and $\Der_2=1$.
Subtracting~\eqref{eq:third} from~\eqref{eq:next} gives
\[
 \binom n2(n-3)!=\frac{n(n-1)(n-3)!}{2}
              =\frac{n!}{2(n-2)}=|b(n-3)|.
\]
Subtracting~\eqref{eq:last} from~\eqref{eq:third} gives
\[
 (n-1)!-\binom n2(n-3)!
 =(n-1)(n-3)!\left((n-2)-\frac n2\right)
 =\frac{(n-3)!(n-4)(n-1)}2.
\]
For $n=4$, the derangements have cycle types $(4)$ and $(2,2)$.
There are $3!=6$ of the first type and $4!/(2^2\,2!)=3$ of the second.
Consequently,
\[
 T(4,2)=6,\qquad T(4,3)=9,\qquad T(4,4)=3.
\]
The count $\binom42(4-3)!=6$ would distinguish one of two transpositions
and count each permutation of type $(2,2)$ twice. This is precisely why
the $n\geq5$ argument cannot be extended to $n=4$.

Finally, $|b(0)|=3$, whereas for $n\geq5$ the ratio of the actual adjacent
difference to the right side of~\eqref{eq:original} is
\[
 \frac{n!/[2(n-2)]}{(n-1)!/[2(n-3)]}
 =\frac{n(n-3)}{n-2}>1.
\]
The strict inequality follows from
$n(n-3)-(n-2)=(n-2)^2-2>0$ for $n\geq5$.
This proves every assertion of the theorem.
\end{proof}

At $n=5$ there are $24$ cycles of type $(5)$ and
$\binom52\,2!=20$ permutations of type $(3,2)$. Hence
\[
 T(5,3)=24,\quad T(5,4)=44,\quad T(5,5)=20,
 \qquad 44-24=20\ne6=|b(1)|.
\]
This is an explicit counterexample to the literal conjecture with its original
offset, not to the corrected formula.

\section{The companion row-gcd conjecture}

The COMMENTS section of the same entry conjectures
\begin{equation}\label{eq:gcd}
 \gcd_{1\leq k\leq n}T(n,k)=n-1\qquad(n\geq3).
\end{equation}
Su's Theorem~7.3~\cite{su} proves exactly~\eqref{eq:gcd}, using the raw and
normalized prime-shift congruences of his Lemma~7.2. This is the earliest
proof found. Kallat's later independent proof is Theorem~4.1(b)~\cite{valuation},
the case $r=1$ of his Theorem~5.1, with the proof in Appendix~A.11.4.
His gcd is over columns $2,\ldots,n$; adding the zero first column leaves it
unchanged. The October 11 shift-congruence paper~\cite[p.~27]{shift}
discusses the triangle before these two proofs.

The common-divisor direction has a short independent explanation.
Let $A_k(s)$ count derangements on $s$ labels avoiding $k$-cycles, for $k\geq2$,
with $A_k(0)=1$. Choosing the cycle containing one specified label gives
\begin{equation}\label{eq:recurrence}
 A_k(n)=\sum_{\substack{2\leq\ell\leq n\\\ell\ne k}}
 \binom{n-1}{\ell-1}(\ell-1)!A_k(n-\ell).
\end{equation}
Indeed, choose its other $\ell-1$ labels, arrange that cycle, and complete
the permutation independently on the remaining labels.
For $n\geq2$ every nonempty summand contains the factor $n-1$, because
\[
 \binom{n-1}{\ell-1}(\ell-1)!
 =(n-1)(n-2)\cdots(n-\ell+1).
\]
Therefore $n-1$ divides every $T(n,k)$ for $k\geq2$.
For the reverse divisibility in~\eqref{eq:gcd}, we use Su's cited theorem;
the last three columns alone do not give it. At $n=7$, for example, their
gcd is $\gcd(1350,1854,1134)=18$, whereas the row gcd is $6$.

\section{Checks and verification scope}

For orientation the two differences are displayed below. The $n=4$ row is
the exceptional boundary, not an instance of the $n\geq5$ formulas.
\begin{center}
\begin{tabular}{rrrrrr}
\toprule
$n$&$T(n,n-2)$&$T(n,n-1)$&$T(n,n)$&adjacent difference&terminal difference\\
\midrule
4&6&9&3&3&3\\
5&24&44&20&20&4\\
6&175&265&145&90&30\\
7&1350&1854&1134&504&216\\
\bottomrule
\end{tabular}
\end{center}
The terminal difference in the last column means $T(n,n-2)-T(n,n)$.
Exact enumeration through $n=9$ and recurrence~\eqref{eq:recurrence} through
$n=80$ independently check the displayed values and the stated identities
on those finite ranges. The all-$n$ difference statements are established by
the proof above; the all-$n$ gcd equality is the cited theorem of Su, with
Kallat's later proof providing a second reference.

The coverage of the three conjectures in revision \#70 is as follows.
\begin{center}
\small
\begin{tabular}{p{0.22\linewidth}p{0.69\linewidth}}
\toprule
Source item & Mathematical coverage\\
\midrule
FORMULA, first conjecture & Refuted and corrected by Theorem~\ref{thm:main};
Su, Proposition~7.4 and Remark~7.5.\\
FORMULA, second conjecture & Proved by Theorem~\ref{thm:main};
Su, Proposition~7.4.\\
COMMENTS, row gcd & Prior proof cited: Su, Theorem~7.3;
common divisibility proved above.\\
\bottomrule
\end{tabular}
\end{center}
The recorded formulas for columns $n-1$ and $n$ are equations
\eqref{eq:next} and~\eqref{eq:last}. The cross-referenced A000166,
A038205, and A122974 describe counts used to define the triangle;
A238474 supplies the offset in~\eqref{eq:b}. This note does not address
other statements in those entries or the Sheffer triangle A238363.

At the pinned revision~\cite{lean}, the Lean declaration
\href{https://github.com/the-omega-institute/trureturing/blob/cdafeab0d44da2b6840cfc802535d2f94a146fad/D5/S3/Combinatorics/Derangements/AdjacentDifferenceRefutation.lean}{\texttt{AdjacentDifferenceRefutation.result}}
expresses the negation of the printed universal adjacent formula, with witness
$n=5$. The declaration
\href{https://github.com/the-omega-institute/trureturing/blob/cdafeab0d44da2b6840cfc802535d2f94a146fad/D5/S3/Combinatorics/Derangements/TerminalDifference.lean}{\texttt{TerminalDifference.result}}
expresses the terminal formula for every $n\geq5$, using rational subtraction.
Both use the namespace \texttt{D5.S3.Combinatorics.Derangements}.
The corrected adjacent formula and the long-cycle identity have written
proofs here; the row-gcd equality is attributed above. The finite checks
cover only their stated ranges.

\begin{thebibliography}{9}
\bibitem{triangle} P. Kagey, \emph{A378495: Triangle of derangements with no
$k$-cycles}, The On-Line Encyclopedia of Integer Sequences, created
November 29, 2024; internal revision \#70, October 17, 2026.
\url{https://oeis.org/A378495/internal}.

\bibitem{bsequence} The On-Line Encyclopedia of Integer Sequences,
\emph{A238474}, formula and offset, accessed November 2, 2026.
\url{https://oeis.org/A238474/internal}.

\bibitem{su} L.-H. Su, \emph{Congruences for Restricted-Cycle and
Pell-Weighted Permutations}, Zenodo, version~1, October 18, 2026.
Theorem~7.3 and Proposition~7.4, p.~9; Remark~7.5, p.~10.
\url{https://doi.org/10.5281/zenodo.23445724}.

\bibitem{valuation} A. Kallat, \emph{A $p$-adic Valuation Calculus for
Permutations with Restricted Cycle Lengths}, arXiv:2610.40342v1,
October 29, 2026. Theorem~4.1, footnote~3, Theorem~5.1,
and Appendices~A.11.3--A.11.4.
\url{https://arxiv.org/abs/2610.40342v1}.

\bibitem{shift} A. Kallat, \emph{Universal Shift Congruences for Weighted
Permutation Sequences}, arXiv:2610.15849v1, October 11, 2026, p.~27.
\url{https://arxiv.org/abs/2610.15849v1}.

\bibitem{lean} Trureturing contributors, \emph{AdjacentDifferenceRefutation.lean}
and \emph{TerminalDifference.lean}, revision
\href{https://github.com/the-omega-institute/trureturing/blob/cdafeab0d44da2b6840cfc802535d2f94a146fad/D5/S3/Combinatorics/Derangements/AdjacentDifferenceRefutation.lean}{adjacent-formula declaration}
and
\href{https://github.com/the-omega-institute/trureturing/blob/cdafeab0d44da2b6840cfc802535d2f94a146fad/D5/S3/Combinatorics/Derangements/TerminalDifference.lean}{terminal-formula declaration}.
The declarations are respectively \texttt{AdjacentDifferenceRefutation.result}
and \texttt{TerminalDifference.result} in the namespace
\texttt{D5.S3.Combinatorics.Derangements}.
\end{thebibliography}

\smallskip
\noindent\small Copyright 2026 Haobo Ma and Wenlin Zhang.
Licensed under Creative Commons Attribution 4.0 International (CC BY 4.0).
\url{https://creativecommons.org/licenses/by/4.0/}.
\end{document}
